% PAGES: 213-214

We note that, if the correlation matrix of the random variables $X_1, X_2, \ldots,
X_n$ is given, the covariance matrix of $X_1, X_2, \ldots, X_n$ is not uniquely
determined.

% NOTE: source reads "Let X1, X2 random variables" (no "be") -- reproduced verbatim.
% UNCLEAR: PDF page 213 (folio 197) mixes the two subscript cases within a single
% example, and this is reproduced as printed. Verified at 2400 dpi: the X-vector
% subscript is a bold LOWERCASE x (\Sigma_{\bfx}, \Omega_{\bfx}) while the
% Y-vector subscript is a bold CAPITAL Y (\Sigma_{\bfY}, \Omega_{\bfY}) -- the Y
% glyph sits on the baseline with no descender, in "Omega_x = Omega_Y" and
% "Sigma_x != Sigma_Y" near the foot of the page as well as in the setup above.
% This is an author inconsistency, not an OCR artifact; do NOT normalise the Y
% to lowercase. The same x/Y split recurs on pages 224, 234, 237, 240 and 241.
\begin{examplex}
Let $X_1$, $X_2$ random variables with covariance matrix $\Sigma_{\bfx} =
\begin{pmatrix} 4 & -3 \\ -3 & 9 \end{pmatrix}$,
and let $Y_1$, $Y_2$ random variables with covariance matrix $\Sigma_{\bfY} =
\begin{pmatrix} 1 & -2 \\ -2 & 16 \end{pmatrix}$.

Then, $\sigma(X_1) = \sqrt{\Sigma_{\bfx}(1,1)} = 2$ and $\sigma(X_2) =
\sqrt{\Sigma_{\bfx}(2,2)} = 3$, and, from (7.21), it follows that the correlation
matrix of $X_1$ and $X_2$ is
\[
\Omega_{\bfx} \;=\; \begin{pmatrix} \frac12 & 0 \\ 0 & \frac13 \end{pmatrix}
\begin{pmatrix} 4 & -3 \\ -3 & 9 \end{pmatrix}
\begin{pmatrix} \frac12 & 0 \\ 0 & \frac13 \end{pmatrix}
\;=\; \begin{pmatrix} 1 & -0.5 \\ -0.5 & 1 \end{pmatrix}.
\]
Similarly, $\sigma(Y_1) = \sqrt{\Sigma_{\bfY}(1,1)} = 1$ and $\sigma(Y_2) =
\sqrt{\Sigma_{\bfY}(2,2)} = 4$, and, from (7.21), it follows that the correlation
matrix of $Y_1$ and $Y_2$ is
\[
\Omega_{\bfY} \;=\; \begin{pmatrix} 1 & 0 \\ 0 & \frac14 \end{pmatrix}
\begin{pmatrix} 1 & -2 \\ -2 & 16 \end{pmatrix}
\begin{pmatrix} 1 & 0 \\ 0 & \frac14 \end{pmatrix}
\;=\; \begin{pmatrix} 1 & -0.5 \\ -0.5 & 1 \end{pmatrix}.
\]
Thus, $X_1$, $X_2$ and $Y_1$, $Y_2$ have the same correlation matrix, since
$\Omega_{\bfx} = \Omega_{\bfY}$, although they have different covariance matrices,
since $\Sigma_{\bfx} \neq \Sigma_{\bfY}$.
\end{examplex}

\begin{lemma}
The correlation matrix of random variables with standard deviation equal to $1$ is
equal to the covariance matrix of the random variables.

In other words, let $\Sigma_{\bfx}$ and $\Omega_{\bfx}$ be the covariance and
correlation matrix of the random variables $X_1, X_2, \ldots, X_n$, respectively. If
$\Sigma_{\bfx}(i,i) = 1$ for all $i = 1:n$, then $\Sigma_{\bfx} = \Omega_{\bfx}$.
\end{lemma}

\begin{proof}
Recall from (7.9) that $\Sigma_{\bfx}(i,i) = \var(X_i) = \sigma_i^2$, for all $i =
1:n$. If $\Sigma_{\bfx}(i,i) = 1$, it follows that $\sigma_i = 1$ for all $i = 1:n$.
Then, $D_{\sigma_{\bfx}} = \diag(\sigma_i)_{i=1:n} = I$, and, from (7.16), we obtain
that $\Omega_{\bfx} = \Sigma_{\bfx}$.
\end{proof}

Note that the result of Lemma 7.2 can also be stated as follows: if all the entries
on the main diagonal of the covariance matrix of $n$ random variables are equal to
1, then the covariance matrix and the correlation matrix of the random variables are
equal.

\begin{lemma}
Let $\Sigma_{\bfx}$ be the covariance matrix of $n$ random variables $X_1, X_2,
\ldots, X_n$. Let $C^{(1)} = \left(c_i^{(1)}\right)_{i=1:n}$ and $C^{(2)} =
\left(c_i^{(2)}\right)_{i=1:n}$ be two column vectors of size $n$. Then,
\begin{equation}
\cov\!\left(\sum_{i=1}^n c_i^{(1)}X_i, \sum_{i=1}^n c_i^{(2)}X_i\right)
\;=\; \left(C^{(1)}\right)^t\Sigma_{\bfx}C^{(2)}
\;=\; \left(C^{(2)}\right)^t\Sigma_{\bfx}C^{(1)}.
\end{equation}

In particular, for $C^{(1)} = C^{(2)} = C = (c_i)_{i=1:n}$, it follows from (7.24)
that
\begin{equation}
\var\!\left(\sum_{i=1}^n c_iX_i\right) \;=\; C^t\Sigma_{\bfx}C.
\end{equation}
\end{lemma}

\begin{proof}
Let $\mu_i = \E{X_i}$, $i = 1:n$. Let
\[
Y^{(1)} \;=\; \sum_{i=1}^n c_i^{(1)}X_i; \qquad Y^{(2)} \;=\; \sum_{i=1}^n c_i^{(2)}X_i.
\]
Then,
\begin{align*}
Y^{(1)} - \E{Y^{(1)}} &= \sum_{i=1}^n c_i^{(1)}(X_i-\mu_i); \\
Y^{(2)} - \E{Y^{(2)}} &= \sum_{i=1}^n c_i^{(2)}(X_i-\mu_i).
\end{align*}
Note that (7.24) is equivalent to
\begin{equation}
\cov\!\left(Y^{(1)},Y^{(2)}\right) \;=\; \left(C^{(1)}\right)^t\Sigma_{\bfx}C^{(2)}
\;=\; \left(C^{(2)}\right)^t\Sigma_{\bfx}C^{(1)}.
\end{equation}
By definition,
\begin{align*}
\cov\!\left(\sum_{i=1}^n c_i^{(1)}X_i, \sum_{i=1}^n c_i^{(2)}X_i\right)
&= \cov\!\left(Y^{(1)},Y^{(2)}\right) \\
&= \E{\left(Y^{(1)}-\E{Y^{(1)}}\right)\left(Y^{(2)}-\E{Y^{(2)}}\right)} \\
&= \E{\left(\sum_{i=1}^n c_i^{(1)}(X_i-\mu_i)\right)\left(\sum_{i=1}^n c_i^{(2)}(X_i-\mu_i)\right)} \\
&= \E{\sum_{1\le j,k\le n} c_j^{(1)}c_k^{(2)}(X_j-\mu_j)(X_k-\mu_k)} \\
&= \sum_{1\le j,k\le n} c_j^{(1)}c_k^{(2)}\,\E{(X_j-\mu_j)(X_k-\mu_k)} \\
&= \sum_{1\le j,k\le n} c_j^{(1)}c_k^{(2)}\cov(X_j,X_k) \\
&= \sum_{1\le j,k\le n} c_j^{(1)}c_k^{(2)}\Sigma_{\bfx}(j,k) \\
&= \left(C^{(1)}\right)^t\Sigma_{\bfx}C^{(2)},
\end{align*}
where, for the last equality, we used formula (10.19), i.e.,
\[
y^tAx \;=\; \sum_{1\le j,k\le n} A(j,k)x_ky_j,
\]
for $A = \Sigma_{\bfx}$ and $x = C^{(2)}$, $y = C^{(1)}$.

Thus, the first equality of (7.24) is proven.

To show that the second equality of (7.24) holds true as well, note that
\begin{equation}
\left(\left(C^{(1)}\right)^t\Sigma_{\bfx}C^{(2)}\right)^t
\;=\; \left(C^{(2)}\right)^t\Sigma_{\bfx}^tC^{(1)}
\;=\; \left(C^{(2)}\right)^t\Sigma_{\bfx}C^{(1)},
\end{equation}
% CONTINUES: \begin{proof} of Lemma 7.3 (opened above) is left OPEN here at the end
% of PDF page 214 (folio 198). Eq. (7.27) above is itself complete and closed; only
% the surrounding prose continues -- "since the covariance matrix Sigma_X is
% symmetric, i.e. ..." on PDF page 215 (folio 199) explains why Sigma_X^t = Sigma_X
% was used in (7.27). The next file (sec07_c1_4.tex, PAGES 215-216) picks up with
% that sentence and closes this proof with \end{proof}. Do not open a new
% \begin{proof} there.
